\documentclass[11pt]{article}

\usepackage[a4paper,margin=1in]{geometry}
\usepackage{amsmath,amssymb}
\usepackage{booktabs}
\usepackage{array}
\usepackage{hyperref}
\usepackage{xcolor}
\usepackage{enumitem}

\hypersetup{
  colorlinks=true,
  linkcolor=blue!50!black,
  urlcolor=blue!50!black,
  citecolor=blue!50!black,
}

\newcommand{\code}[1]{\texttt{#1}}

\title{Methodology: Dynamic Bayesian-Network Structural\\ Causal Model (Dynamic BN-SCM)}
\author{}
\date{}

\begin{document}
\maketitle

\noindent
This section documents the methodology \emph{as implemented} in the
\code{BayesianNetwork-SCM/bayesian\_network\_scm/} package and exercised by
\code{run\_dynamic\_bn\_scm.py} and \code{run\_hypothesis\_experiments.py}.
All numbers and defaults below are taken directly from the source code.

\section{Design overview}

The model is a leakage-controlled, longitudinal structural causal model for
multimodal Alzheimer's disease progression. The primary object is \emph{not}
a same-time cross-sectional DAG, which is known to be vulnerable to reverse
orientation and unmeasured confounding. Instead, for each subject we form a
\textbf{baseline-to-follow-up pair} and model
\[
\text{baseline state } X(t_0) \;\longrightarrow\; \text{annualized rate } R = \frac{Y(t_1) - Y(t_0)}{\Delta t},
\]
with all candidate parents pinned to the baseline visit $t_0$. Future target
values can therefore never enter the predictor set, which gives the graph
an unambiguous temporal orientation and prevents target leakage.

For an outcome $j$ (a regional tau rate or a cognitive rate) and a
subject-pair $i$, the structural equation is
\begin{equation}
R_{ij} \;=\; a_j(Z_i) \;+\; c_j(Z_i)\, Y_{ij}(t_0) \;+\; \sum_{\ell \in \mathrm{pa}(j)} \gamma_{j\ell}\, b_{j\ell}(Z_i)\, X_{i\ell}(t_0) \;+\; \varepsilon_{ij},
\label{eq:scm}
\end{equation}
where
\begin{itemize}[leftmargin=*,nosep]
  \item $R_{ij} = (Y_{ij}(t_1) - Y_{ij}(t_0))/\Delta t_i$ is the annualized rate target,
  \item $Z_i \in [0,1]$ is a train-only \textbf{explainable disease pseudotime} (\S\ref{sec:pseudotime}),
  \item $a_j(Z)$ is a disease-stage-varying intercept (trajectory),
  \item $c_j(Z)\, Y_{ij}(t_0)$ is a pseudotime-varying \textbf{self-history / autonomous} progression term,
  \item $b_{j\ell}(Z)$ is the pseudotime-varying direct effect from baseline parent $\ell$ on target $j$, restricted to the biologically admissible parent set $\mathrm{pa}(j)$,
  \item $\gamma_{j\ell}$ is an edge-inclusion proxy obtained by thresholding the bootstrap stability of $|b_{j\ell}(Z)|$ over $Z$ (\S\ref{sec:stability}).
\end{itemize}
The corresponding biomarker forecast is the explicit Euler step
\begin{equation}
Y_{ij}(t_1) \;=\; Y_{ij}(t_0) \;+\; \Delta t_i \cdot R_{ij}.
\label{eq:forecast}
\end{equation}
Equation~\eqref{eq:scm} is a varying-coefficient extension of a standard
linear SCM: fixing $Z$ recovers an ordinary linear regression of rate on
baseline state, whereas allowing the coefficients to vary with $Z$ lets the
same causal parents have stage-dependent effect magnitudes (e.g.,
early-only or late-only mechanisms).

\section{Data construction}

The pair table is assembled in \code{data.py} by
\code{build\_multimodal\_pair\_dataset()} from the ENIGMA-format ADNI cohort
exports under \code{experiments/group\_average\_enigma/output/}.

\paragraph{Pairs.} Each row is one (baseline, follow-up) tau-PET pair. We
require \code{baseline\_loniuid} and \code{target\_loniuid} to resolve in
\code{cohort\_tau\_observations.csv}, finite baseline and follow-up dates,
and $\Delta t_i = \texttt{target\_time\_years} > 0$. Tau follow-up is always
present because the cohort was constructed from longitudinal FTP tau-PET
pairs.

\paragraph{Baseline predictors.} For every pair $i$, baseline predictors are
taken from the visit nearest in time to the baseline tau-PET date, within
\code{max\_date\_distance\_days = 1095} (3~years).

\begin{table}[h]
\centering
\small
\begin{tabular}{>{\raggedright}p{0.20\linewidth} p{0.72\linewidth}}
\toprule
\textbf{Layer} & \textbf{Features (30 total)} \\
\midrule
root & \code{age\_years}, \code{sex\_female}, \code{education\_years}, \code{apoe4\_dose} \\
fluid & \code{plasma\_pt217}, \code{plasma\_ab42\_ab40}, \code{plasma\_nfl}, \code{plasma\_gfap} \\
pathology & \code{amyloid\_summary\_suvr}, \code{amyloid\_centiloids}, \code{amyloid\_positive}, \code{tau\_meta\_temporal}, plus 10 regional tau SUVR \\
neurodegeneration & \code{mri\_hippocampus\_volume}, \code{mri\_amygdala\_volume}, \code{mri\_temporal\_cortical\_volume}, \code{mri\_temporal\_cortical\_thickness} \\
clinical & \code{adas13}, \code{mmse}, \code{ravlt\_immediate}, \code{cdrsb} \\
\bottomrule
\end{tabular}
\end{table}

\code{apoe4\_dose} is the count of $\varepsilon 4$ alleles in the APOE
genotype. MRI volumes and thicknesses are means over left/right entorhinal,
fusiform, inferiortemporal, and middletemporal cortical labels parsed from
the FreeSurfer dictionary; hippocampal and amygdalar segmentation volumes
are bilateral averages. Rows with overall FreeSurfer QC \code{FAIL} and
amyloid scans with QC fail flags are excluded.

\paragraph{Regional tau predictors and targets.} Ten temporo-parietal aparc
labels are used: \code{L/R\_entorhinal}, \code{L/R\_fusiform},
\code{L/R\_inferiortemporal}, \code{L/R\_middletemporal},
\code{L/R\_inferiorparietal}, mapped to ADNI tau columns via
\code{adni\_to\_enigma\_aparc\_mapping.csv}. The target vector contains 15
entries: \code{tau\_rate:meta\_temporal} (composite SUVR rate),
\code{tau\_rate:\textit{region}} for each of the 10 regions above, and
\code{cognitive\_rate:\{adas13, mmse, ravlt\_immediate, cdrsb\}}. Each
target is the simple annualized rate $R_{ij} = (Y_{ij}(t_1) -
Y_{ij}(t_0))/\Delta t_i$. Cognitive rates are computed only when both
nearest-to-baseline and nearest-to-follow-up scores are finite within the
1095-day window.

\paragraph{Subject-level split.} \code{make\_subject\_split()} partitions the
\textbf{unique RIDs} (not pairs) into 60\,\% train / 20\,\% validation /
20\,\% test with random seed \code{20260519}. All pairs from a given
subject are kept inside a single fold, which is essential for leakage
control: the same subject never appears as both a training row and a
held-out row.

\section{Explainable pseudotime \texorpdfstring{$Z$}{Z}}
\label{sec:pseudotime}

\code{fit\_pseudotime()} in \code{pseudotime.py} defines $Z$ so the model
can capture stage-dependent dynamics without using calendar time or future
information.

\paragraph{Feature selection.} Given a \emph{mode}
(\code{tau\_free}, \code{global}, \code{clinical\_free}, \code{pt217\_free}),
only features admissible for the mode with training finite-value coverage
$\geq 0.5$ and non-zero training variance are retained. \code{tau\_free}
(the default) excludes all \code{tau}, \code{pT217}, and \code{pTau}
features so that $Z$ is not endogenous to the tau targets of interest.

\paragraph{Standardization and imputation.} Missing entries are filled with
the column-wise training median; columns are then mean-centered and scaled
to unit standard deviation (with a guard $\text{scale} > 10^{-12}$).

\paragraph{Linear pseudotime via SVD.} Let
$\tilde X_{\text{train}} \in \mathbb{R}^{n_{\text{train}} \times p}$ be the
standardized training feature matrix. We take the leading right singular
vector $w$ of $\tilde X_{\text{train}}$:
\begin{equation}
\tilde X_{\text{train}} \;=\; U\,\Sigma\,V^\top, \qquad w = V_{:,1}.
\label{eq:svd}
\end{equation}
For any (training or held-out) row, the raw score is the projection
$s_i = \tilde X_i^\top w$. Scores are rescaled to $[0,1]$ using the
empirical 1st and 99th training percentiles $(s_{\text{lo}}, s_{\text{hi}})$:
\begin{equation}
Z_i \;=\; \mathrm{clip}\!\left(\frac{s_i - s_{\text{lo}}}{s_{\text{hi}} - s_{\text{lo}}},\,0,\,1\right).
\label{eq:zscore}
\end{equation}
This is the SVD-based, linear analog of a single principal axis of disease
progression; using only the leading component makes $Z$ one-dimensional and
interpretable.

\paragraph{Orientation.} Because singular vectors are sign-ambiguous, $w$
is flipped so its correlation with a sign-coded disease-burden proxy is
non-negative. The burden proxy is the mean of $z$-scored features signed by
\code{burden\_sign()} (positive: \code{tau}, \code{pT217}, \code{amyloid},
\code{centiloid}, \code{nfl}, \code{gfap}, \code{adas13}, \code{cdrsb};
negative: \code{mmse}, \code{ravlt}, \code{volume}, \code{thickness},
\code{ab42\_ab40}). This guarantees that higher $Z$ corresponds to higher
disease burden, while leaving the geometry of $Z$ itself unchanged.

\paragraph{Reportable diagnostics.} The model exposes feature loadings $w$,
per-subject contributions $w_m \cdot z(X_{im})$, a diagnosis-group ordering
of median $Z$ (CN $<$ MCI $<$ AD is the expected sanity check), and the
explained-variance ratio of the leading singular value.

\section{Biological graph constraints}

\code{CausalConstraints} in \code{constraints.py} restricts the set of
admissible baseline parents for each rate target.

\paragraph{Variable layers.} Every variable is assigned a layer with rank
$\text{root}(0) < \text{fluid}(1) < \text{pathology}(2) < \text{neurodegeneration}(3) < \text{clinical}(4)$.
Layer inference uses keyword rules
(e.g.\ \code{plasma}/\code{pt217}/\code{nfl}/\code{gfap} $\to$ \code{fluid};
\code{amyloid}/\code{tau} $\to$ \code{pathology};
\code{mri}/\code{volume}/\code{thickness} $\to$ \code{neurodegeneration};
\code{adas}/\code{mmse}/\code{ravlt}/\code{cdrsb} $\to$ \code{clinical}).

\paragraph{Hard constraints.} For \code{can\_parent(parent, child)}, parent
$\ell$ may direct an edge to child $j$ only if:
\begin{enumerate}[leftmargin=*,nosep]
  \item $\text{parent} \neq \text{child}$;
  \item parent is not a clinical sink and child is not a root;
  \item either $\mathrm{layer}(\text{parent}) < \mathrm{layer}(\text{child})$ (cross-layer downstream), or
  \item parent and child are in the same layer \emph{and} the same-layer rule below is satisfied.
\end{enumerate}

\paragraph{Same-layer rules.}
\begin{itemize}[leftmargin=*,nosep]
  \item \code{root}, \code{clinical}, \code{neurodegeneration}: no within-layer edges.
  \item \code{fluid}: ordered as \code{amyloid/ab42} $<$ \code{pT217/pTau} $<$ \code{gfap} $<$ \code{nfl}; only parent-to-later-stage child edges are allowed.
  \item \code{pathology}: only \code{amyloid} $\to$ \code{tau} is allowed; tau cannot drive amyloid; no edges within amyloid or within tau.
\end{itemize}
This encodes the canonical AD cascade
$\text{root} \to \text{fluid} \to \text{amyloid} \to \text{tau} \to \text{neurodegeneration} \to \text{clinical}$,
with diagnosis left out of the parent set entirely and reserved for
validation and stratification.

\section{Varying-coefficient design matrix}

\code{build\_design\_matrix()} in \code{dynamic\_scm.py} constructs the
per-target design that operationalizes Eq.~\eqref{eq:scm}.

\paragraph{B-spline basis on $Z$.} Let $B(Z) \in \mathbb{R}^{n \times K}$
be a cubic B-spline basis with $K$ columns, built by
\code{sklearn.preprocessing.SplineTransformer} with $\text{degree}=3$,
$\text{n\_knots}=4$, $\text{include\_bias}=\text{True}$, and constant
extrapolation. The varying intercept and the varying parent effects are
expanded as
\begin{equation}
a_j(Z) = \sum_{k=1}^{K} \alpha_{jk} B_k(Z), \quad
c_j(Z) = \sum_{k=1}^{K} \kappa_{jk} B_k(Z), \quad
b_{j\ell}(Z) = \sum_{k=1}^{K} \theta_{j\ell k} B_k(Z).
\label{eq:basis}
\end{equation}
If the training set has too few unique $Z$ values to support the basis, the
implementation falls back to a polynomial basis of width
$\min(4,\,\text{degree}+1)$.

\paragraph{Design columns.} For each target $j$ with chosen parent set
$\mathrm{pa}(j)$, the design row $x_i$ is the concatenation of three blocks:
\begin{enumerate}[leftmargin=*,nosep]
  \item \textbf{Trajectory:} the $K$ basis values $B_k(Z_i)$ $\Rightarrow a_j(Z)$.
  \item \textbf{Self-history:} $Y_{ij}(t_0)\, B_k(Z_i)$ for $k=1,\dots,K$ $\Rightarrow c_j(Z)\,Y_{ij}(t_0)$.
  \item \textbf{Edge-by-stage interactions:} $X_{i\ell}(t_0)\, B_k(Z_i)$ for each parent $\ell \in \mathrm{pa}(j)$ and $k=1,\dots,K$ $\Rightarrow b_{j\ell}(Z)\,X_{i\ell}(t_0)$.
\end{enumerate}
The total design dimension for target $j$ is therefore
$D_j = K\,(2 + |\mathrm{pa}(j)|)$. The recovered effect curves
(\S\ref{sec:stability}) are computed by contracting these coefficients back
through $B(Z_{\text{grid}})$.

\section{Parent selection and ridge estimation}

\paragraph{Candidate parents.} For target $j$, \code{candidate\_parents()}
returns every feature $\ell$ that satisfies the layer/role constraints of
\S4 \emph{and} has training coverage $\geq 0.25$. Features are then ranked
by the absolute Pearson correlation between $X_\ell(t_0)$ and $R_j$ on the
training rows. A small set of biological priority features is prepended
ahead of the correlation ranking (for \code{tau\_rate:*}:
\code{plasma\_pt217}, \code{amyloid\_summary\_suvr},
\code{amyloid\_centiloids}, \code{plasma\_ab42\_ab40}, \code{apoe4\_dose},
\code{age\_years}; for \code{cognitive\_rate:*}: \code{tau\_meta\_temporal},
\code{mri\_hippocampus\_volume}, \code{amyloid\_summary\_suvr},
\code{plasma\_nfl}, \code{age\_years}). The final list is truncated to
\code{max\_parents\_per\_target} (default $8$ in
\code{run\_dynamic\_bn\_scm.py}, $6$ in
\code{run\_hypothesis\_experiments.py}).

\paragraph{Preprocessing.} Within \code{fit\_ridge()}, rows with non-finite
targets are dropped, design entries are imputed with column-wise training
medians, and columns are mean-centered and scaled by their training
standard deviation (with a $> 10^{-12}$ guard).

\paragraph{Ridge solution.} Let
$\Phi_{\text{train}} \in \mathbb{R}^{n \times D}$ be the standardized
training design and $r_{\text{train}}$ the centered training rate target.
The estimator solves
\begin{equation}
\hat{\beta} \;=\; \arg\min_{\beta} \;\bigl\lVert r_{\text{train}} - \Phi_{\text{train}}\,\beta \bigr\rVert^2 + \alpha\,\lVert \beta \rVert^2 \;=\; \bigl(\Phi_{\text{train}}^\top \Phi_{\text{train}} + \alpha\, I_D\bigr)^{-1} \Phi_{\text{train}}^\top r_{\text{train}},
\label{eq:ridge}
\end{equation}
with intercept $\beta_0 = \mathrm{mean}(r_{\text{train}})$. Ridge
regularization is used in place of an explicit second-order random-walk
prior because the prototype estimates effects pointwise rather than
sampling a posterior; the $L_2$ penalty plays the role of the
inverse-variance hyperprior on the spline coefficients while remaining
convex and closed form.

\paragraph{Penalty selection.} The penalty is chosen by
\textbf{subject-grouped 5-fold cross-validation} using
\code{sklearn.model\_selection.GroupKFold} with groups equal to each pair's
\code{RID}. This guarantees that no subject appears in both the inner train
and inner validation folds. The grid is
$\alpha \in \{1, 10, 100, 1000, 10000\}$ and the $\alpha$ minimizing
mean-squared prediction error across folds is retained.

\paragraph{Effect-curve reconstruction.} Once $\hat{\beta}$ is fit, the
effect of parent $\ell$ at any $Z$ is recovered as
\begin{equation}
b_{j\ell}(Z) \;=\; \sum_{k=1}^{K} B_k(Z) \cdot \frac{\hat{\beta}[\mathrm{col}(\ell,k)]}{\mathrm{scale}[\mathrm{col}(\ell,k)]},
\label{eq:effect}
\end{equation}
with the analogous formula for $c_j(Z)$. The grid
$Z = \mathrm{linspace}(0,1,101)$ is used to evaluate all curves.

\section{Edge stability and inclusion (\texorpdfstring{$\gamma_{j\ell}$}{gamma})}
\label{sec:stability}

A first inclusion screen is the \textbf{effect-magnitude threshold}
$\tau_{\text{eff}}$: edge $\ell \to j$ is ``selected by effect'' if
$\max_Z |b_{j\ell}(Z)| \geq \tau_{\text{eff}}$, with $\tau_{\text{eff}} = 0.01$
by default.

A frequentist analog of an edge posterior inclusion probability (PIP) is
then obtained by \textbf{subject-block bootstrap}
(\code{bootstrap\_edge\_stability()}):
\begin{enumerate}[leftmargin=*,nosep]
  \item Group training pairs by \code{RID}; let $\mathcal{G}$ be the set of training subjects.
  \item For $b = 1,\dots,B$ iterations (default $B = 25$), sample $|\mathcal{G}|$ subjects from $\mathcal{G}$ with replacement, take all of each sampled subject's training pairs to form a bootstrap training set, and re-fit the full Dynamic BN-SCM on this sample (with \code{cv\_folds = 0} for speed).
  \item For each candidate parent--target pair $(\ell, j)$, record whether $\max_Z |b_{j\ell}^{(b)}(Z)| \geq \tau_{\text{eff}}$.
\end{enumerate}
The empirical \textbf{bootstrap inclusion probability} is
\begin{equation}
\mathrm{PIP}^{\text{boot}}_{j\ell} \;=\; \frac{1}{B} \sum_{b=1}^{B} \mathbf{1}\!\left\{ \max_Z \bigl| b_{j\ell}^{(b)}(Z) \bigr| \geq \tau_{\text{eff}} \right\},
\label{eq:pip}
\end{equation}
and is used as the implementation's proxy for $\gamma_{j\ell}$. Bootstrap
blocks are at the subject level so resampled pairs respect the
subject-level dependence structure of the cohort.

\section{Forecasting and evaluation}

For any held-out pair, the predicted SUVR (or score) follow-up is obtained
from Eq.~\eqref{eq:forecast} using the model's predicted rate
$\hat{R}_{ij}$ together with the observed $\Delta t_i$ and $Y_{ij}(t_0)$.
Evaluation metrics are computed per (target, split) pane and aggregated
across targets:
\begin{itemize}[leftmargin=*,nosep]
  \item per-pair rate error metrics \code{mae}, \code{rmse}, \code{pearson}, \code{spearman} between observed and predicted rates;
  \item per-subject regional \code{delta\_spearman} between observed and predicted baseline-to-follow-up change vectors;
  \item a compact summary of validation / test tau-rate MAE and Spearman (mean and median across targets) for ranking and model comparison.
\end{itemize}

\section{Pre-specified hypothesis tests}

\code{run\_hypothesis\_experiments.py} runs three pre-registered analyses on
top of the base \code{tau\_free} fit.

\subsection{H1 --- pT217-to-tau decoupling}

For every tau-rate target $j$ and the parent $\ell = $ \code{plasma\_pT217},
the effect curve $b_{\text{pT217},j}(Z)$ is sliced into three windows
$\{$early: $Z \leq 0.30$;\, mid: $0.30 < Z < 0.70$;\, late: $Z \geq 0.70\}$.
We report
\begin{itemize}[leftmargin=*,nosep]
  \item the means of $|b_{\text{pT217},j}(Z)|$ in each window;
  \item the indicator $\text{late\_less\_than\_early} = (\text{mean}_{\text{late}} < \text{mean}_{\text{early}})$;
  \item a ``near-zero late'' indicator $\text{mean}_{\text{late}} < 0.005$;
  \item a strict decoupling pattern that additionally requires the self-history effect to rise at late $Z$;
  \item the decoupling coordinate
\begin{equation}
Z_{\text{decouple}} \;=\; \min\bigl\{ Z \geq 0.30 \;:\; |b_{\text{pT217},j}(Z)| < 0.005 \bigr\}.
\label{eq:zdecouple}
\end{equation}
\end{itemize}
The prototype reports these as descriptive summaries of the ridge effect
curves; turning them into a formal posterior probability would require a
Bayesian sampler over $\theta_{j\ell}$, which is explicitly out of scope
for the current implementation.

\subsection{H2 --- Transcriptomic resilience gating}

The intended structural extension introduces a regional gene-expression
modifier $SR_r$ of the amyloid $\to$ tau edge:
\begin{equation}
R_{ij,r} \;=\; \cdots \;+\; b_{A\beta,\text{tau},r}(Z_i)\,\text{amyloid}_i(t_0) \;+\; d_{SR,r}(Z_i)\,\text{amyloid}_i(t_0)\,SR_r \;+\; \cdots,
\label{eq:h2}
\end{equation}
with expected sign $d_{SR,r}(Z) < 0$ (regions with higher resilience
expression have lower amyloid-driven tau acceleration).
\code{evaluate\_h2\_data\_gate()} scans the project tree for
AHBA~/~Allen~/~gene-expression files and reports H2 as
\code{not\_testable\_current\_data} if none are found. The code
intentionally refuses to fit a structural-proxy H2 model until a
region-by-DK/aparc expression matrix is available.

\subsection{H3 --- PART-like vs.\ AD-continuum spatial route}

Each pair is assigned an A/T status using the implemented data:
\begin{itemize}[leftmargin=*,nosep]
  \item $A+$ if the baseline \code{amyloid\_positive} feature $\geq 0.5$, else $A-$;
  \item $T+$ if baseline \code{tau\_meta\_temporal} $\geq q_{75}$ of the \textbf{training cognitively-normal subset}, else $T-$.
\end{itemize}
Subjects with non-finite amyloid or tau are labelled \code{unclassified}.
The two analyzed groups are $A{-}T{+}$ (PART-like) and $A{+}T{+}$ (AD
continuum). For each group, regional rate vectors are computed by averaging
both the observed and the BN-SCM-predicted regional rates across pairs.
Two quantities are reported:
\begin{align}
\text{route\_similarity} &\;=\; \rho_{\text{Spearman}}\bigl( \overline{R}_{\text{PART}},\, \overline{R}_{\text{AD}} \bigr), \label{eq:h3route} \\
\text{kinetic\_ratio}_{\text{AD/PART}} &\;=\; \frac{\mathrm{rms}(\overline{R}_{\text{AD}})}{\mathrm{rms}(\overline{R}_{\text{PART}})}. \label{eq:h3kin}
\end{align}
H3 is declared ``supported'' only when both
$\text{route\_similarity} \geq 0.70$ and
$\text{kinetic\_ratio} \geq 1.20$ on the test split, encoding the
``shared spatial route with kinetic acceleration'' claim.
Equation~\eqref{eq:h3route} is computed both for observed-rate route
similarity and for the BN-SCM-predicted-rate route, so that the model's
spatial inductive bias can be checked against the data.

\section{Sensitivity analyses}

Two robustness sweeps are run alongside H1--H3.

\paragraph{Pseudotime sensitivity.} The full Dynamic BN-SCM is re-fit four
times, once per pseudotime mode (\code{tau\_free}, \code{global},
\code{clinical\_free}, \code{pt217\_free}), reporting the median test rate
MAE/RMSE and the median test delta-Spearman. This isolates the effect of
$Z$ from the effect of the parent set.

\paragraph{Parent ablations.} Starting from the \code{tau\_free} model,
candidate features are deleted in groups (\code{\{plasma\_pt217\}},
\code{\{amyloid\_summary\_suvr, amyloid\_centiloids, amyloid\_positive\}},
\code{\{plasma\_ab42\_ab40\}}, \code{\{apoe4\_dose\}}) and the model is
re-fit and re-evaluated. Differences in test rate MAE versus the full
model quantify how much each modality contributes to held-out tau-rate
forecasting.

\section{Leakage controls and limitations}

By construction, the implementation enforces:
\begin{itemize}[leftmargin=*,nosep]
  \item predictors are sampled at or before $t_0$, never at $t_1$;
  \item the subject split partitions unique \code{RID}s;
  \item pseudotime, ridge centering/scaling, $\alpha$ selection, parent ranking, and bootstrap blocks are all fit using \textbf{training rows only};
  \item ridge $\alpha$ selection uses \code{GroupKFold} with \code{RID} groups so no subject's pairs straddle the inner CV boundary;
  \item diagnosis is used only for stratification and reporting, never as a parent.
\end{itemize}

Stated limitations carried directly in the code are:
\begin{enumerate}[leftmargin=*,nosep]
  \item The estimator is a penalized (ridge) varying-coefficient regression with subject-block bootstrap stability, not a full posterior MCMC over the graph; reported edge probabilities should be read as bootstrap inclusion proxies, not Bayesian PIPs.
  \item H2 is gated off until a DK/aparc-aligned regional gene-expression matrix is provided.
  \item Brain-map panels visualize only the ten temporo-parietal target regions; un-modelled DK regions are greyed out.
  \item The H3 $T+$ threshold is currently derived from the training CN 75th percentile of \code{tau\_meta\_temporal}; a project-level tau-positivity threshold would be substituted if configured.
\end{enumerate}

\end{document}
